\documentclass{article}
\usepackage{arxiv}
%\usepackage{amsthm}
\usepackage{amssymb}
\usepackage[utf8]{inputenc} % allow utf-8 input
\usepackage[T1]{fontenc}    % use 8-bit T1 fonts
\usepackage{mathtools}
\usepackage{hyperref}  
\usepackage{mathrsfs}

\usepackage{graphicx}
\usepackage{subcaption}
% hyperlinks
\usepackage{url}            % simple URL typesetting
\usepackage{booktabs}       % professional-quality tables
\usepackage{amsfonts}       % blackboard math symbols
\usepackage{nicefrac}       % compact symbols for 1/2, etc.
\usepackage{microtype}      % microtypography
\usepackage{lipsum}		% Can be removed after putting your text content
\usepackage{graphicx}
\usepackage{natbib}
\usepackage{doi}
\usepackage{amsmath}
\usepackage{physics}
\usepackage{wrapfig}
\usepackage{bm}
\usepackage{algorithm,algpseudocode}
\usepackage{xcolor}
\usepackage{forest}
\usepackage{tcolorbox}
\usepackage{tikz-cd}
\usepackage{enumitem}


\newenvironment{proof}{\paragraph{Proof:}}{\hfill$\square$}
\newenvironment{solution}{\paragraph{Solution:}}{\hfill$\square$}


%—– Blackboard‐bold sets —–
\newcommand{\R}{\mathbb{R}}
\newcommand{\N}{\mathbb{N}}
\newcommand{\Z}{\mathbb{Z}}
\newcommand{\Q}{\mathbb{Q}}
\newcommand{\C}{\mathbb{C}}

%—– Probability & expectation —–
\newcommand{\Prb}{\mathbb{P}}        % probability
\newcommand{\E}{\mathbb{E}}           % expectation
\newcommand{\Var}{\mathrm{Var}}       % variance
\newcommand{\Cov}{\mathrm{Cov}}       % covariance

%—– Common math operators —–
\newcommand{\diag}{\operatorname{diag}}
\newcommand{\supp}{\operatorname{supp}}
\newcommand{\argmax}{\operatorname*{arg\,max}}
\newcommand{\argmin}{\operatorname*{arg\,min}}
\newcommand{\ind}{\mathbf{1}}         % indicator

%—– Norms, inner products, absolute value —–

\newcommand{\inner}[2]{\left\langle #1, #2 \right\rangle}

%—– Calligraphic sets —–
\newcommand{\cA}{\mathcal{A}}
\newcommand{\cS}{\mathcal{S}}
\newcommand{\cT}{\mathcal{T}}
\newcommand{\cF}{\mathcal{F}}
\newcommand{\cE}{\mathcal{E}}

%—– RL / MDP notation —–
\newcommand{\SSS}{\mathcal{S}}        % state space
\newcommand{\AAA}{\mathcal{A}}        % action space
\newcommand{\PP}{\mathcal{P}}         % transition kernel
\newcommand{\RRR}{\mathcal{R}}        % reward function
\newcommand{\gammaRL}{\gamma}         % discount factor

%—– Confidence bounds —–
\newcommand{\olQ}{\overline{Q}}
\newcommand{\ulQ}{\underline{Q}}
\newcommand{\olV}{\overline{V}}
\newcommand{\ulV}{\underline{V}}

%—– Miscellaneous —–
\newcommand{\eps}{\varepsilon}
\newcommand{\dt}{\mathrm{d}t}
\newcommand{\given}{\,\bigl\vert\,}    % conditioning bar

%—– Shorthand for frequently used fractions —–
\newcommand{\half}{\tfrac12}
\newcommand{\thalf}{\tfrac32}

%—– Styling —–
\newcommand{\todo}[1]{\textcolor{red}{[TODO: #1]}}


\newcommand{\note}[1]{\begin{tcolorbox}
\textbf{Note:} #1\end{tcolorbox}}

\newtheorem{theorem}{Theorem}
\newtheorem{proposition}[theorem]{Proposition}
\newtheorem{lemma}[theorem]{Lemma}
\newtheorem{corollary}[theorem]{Corollary}
\newtheorem{definition}[theorem]{Definition}
\newtheorem{assumption}[theorem]{Assumption}
\newtheorem{remark}[theorem]{Remark}
\newtheorem{example}[theorem]{Example}
\newtheorem{conjecture}[theorem]{Conjecture}
\title{A title}
\date{} 					% Or removing it

\author{ \href{https://orcid.org/0000-0003-1700-756X}{William Chang} \\
	Department of Applied Mathematics\\
	University of California, Los Angeles\\
	Los Angeles, CA \\
	\url{https://williamc.me/} \\
	%% examples of more authors
	\And
}

% Uncomment to override  the `A preprint' in the header
\renewcommand{\headeright}{}
\renewcommand{\undertitle}{}
%\renewcommand{\shorttitle}{\textit{arXiv} Template}

%%% Add PDF metadata to help others organize their library
%%% Once the PDF is generated, you can check the metadata with
%%% $ pdfinfo template.pdf
\hypersetup{
pdftitle={},
pdfsubject={},
pdfauthor={William Chang},
pdfkeywords={},
}

\begin{document}
\maketitle
\begin{abstract}
    
\end{abstract}

% ============================================================
\section{Motivation}
% ============================================================

Reinforcement learning in continuous state--action spaces is fundamentally challenging
because classical tabular methods do not scale to uncountably large spaces. A natural
response is to impose metric structure and exploit regularity of the value function.
\citet{cao2021} formalize this by working in metric spaces $(\mathcal{S} \times
\mathcal{A}, D)$ and requiring the optimal $Q^*$-function to be Lipschitz:

\begin{assumption}[Lipschitz $Q^*$, \citet{cao2021} Assumption 2]
\label{asm:lipschitz}
For every stage $h \in [H]$, $Q_h^*$ is $L$-Lipschitz with respect to $D$:
\[
  \bigl|Q_h^*(x,a) - Q_h^*(x',a')\bigr| \;\le\; L \cdot D\bigl((x,a),(x',a')\bigr).
\]
\end{assumption}

Under this structure the key complexity measure is not the covering dimension of the
full space $\mathcal{S} \times \mathcal{A}$, but the \emph{zooming dimension}, which
captures how rapidly the near-optimal region grows as the sub-optimality threshold is
relaxed.

\begin{definition}[Near-optimal set and zooming dimension, \citet{cao2021} Definitions 3--4]
\label{def:zooming}
The gap function is $\mathrm{gap}_h(x,a) := V_h^*(x) - Q_h^*(x,a)$.
The near-optimal set at scale $r$ is
\[
  P^{Q^*}_{h,r} \;:=\;
    \Bigl\{(x,a)\in\mathcal{S}\times\mathcal{A} :
      \mathrm{gap}_h(x,a) \;\le\;
      \Bigl(\tfrac{2(H+1)}{d_{\max}}+2L\Bigr)r\Bigr\}.
\]
The \emph{zooming dimension} $z_c$ is the infimum over $d>0$ such that
$N^{\mathrm{pack}}_r(P^{Q^*}_{h,r}) \le c\,r^{-d}$ for all $r\in(0,d_{\max}]$.
\end{definition}

Because $P^{Q^*}_{h,r}\subseteq \mathcal{S}\times\mathcal{A}$, the zooming dimension
satisfies $z_c \le d_c$ (the covering dimension), often strictly. This yields an
instance-dependent adaptive regret guarantee:

\begin{corollary}[\citet{cao2021} Corollary 2]
\label{cor:zooming_regret}
For any constant $c>0$, Adaptive Q-learning achieves
\[
  \mathrm{Reg}(K) \;\le\;
  \tilde{\mathcal{O}}\!\left(H^{5/2} + \Lambda^{1/2}\,
    K^{\frac{z_c+1}{z_c+2}}\right),
\]
a polynomial improvement over the non-adaptive rate $K^{(d_c+1)/(d_c+2)}$ whenever
$z_c < d_c$.
\end{corollary}

On the representation side, raw kinematic or visual observations are typically far
richer than the true task-relevant state. \citet{zhang2021} propose \emph{Deep
Bisimulation for Control} (DBC), which learns an encoder $\varphi:\mathcal{S}\to
\mathcal{Z}$ such that $\ell_1$-distances in latent space match bisimulation
distances. The bisimulation metric $\tilde{d}$ is defined by the fixed-point equation:

\begin{definition}[Bisimulation metric, \citet{zhang2021} Definition 2]
\label{def:bisim}
\[
  \tilde{d}(s_i,s_j) \;=\; \max_{a\in\mathcal{A}}\,
    \bigl[(1-c)\,|R^a_{s_i} - R^a_{s_j}|
    \;+\; c\,W_1(P^a_{s_i}, P^a_{s_j};\,\tilde{d})\bigr],
\]
where $W_1$ denotes the $1$-Wasserstein metric.
\end{definition}

The central property linking bisimulation to zooming is that $V^*$ is Lipschitz with
respect to $\tilde{d}$:

\begin{theorem}[\citet{zhang2021} Theorem 5]
\label{thm:vlip}
If $c \ge \gamma$, then $V^*$ is Lipschitz with respect to $\tilde{d}$ with constant
$\frac{1}{1-c}$:
\[
  |V^*(s_i) - V^*(s_j)| \;\le\; \frac{1}{1-c}\,\tilde{d}(s_i, s_j).
\]
\end{theorem}

Theorem~\ref{thm:vlip} directly licenses using the bisimulation metric as the metric
$D$ in Assumption~\ref{asm:lipschitz}: states close in $\tilde{d}$ have close optimal
values, so a single zooming cell can be assigned one optimistic Q-estimate without
large approximation error. Furthermore, \citet{zhang2021} show that the bisimulation
partition corresponds to the causal feature set of the reward:

\begin{theorem}[\citet{zhang2021} Theorem 3]
\label{thm:causal}
If we partition observations using the bisimulation metric, those clusters correspond
to the causal feature set of the observation space with respect to current and future
reward, i.e.\ the minimal sufficient statistic of $V^*$.
\end{theorem}

Theorem~\ref{thm:causal} confirms that the encoder discards task-irrelevant variation
while retaining everything causally upstream of the reward, which is precisely the
compression needed to make joint-space zooming tractable. The combination of these
two results motivates a three-phase pipeline: (1) learn a bisimulation-invariant
low-dimensional representation; (2) cluster the latent space to form a coarse state
discretization; and (3) run an optimistic adaptive zooming algorithm on the continuous
action dimension within each cluster.

% ============================================================
\section{Initial Problems}
% ============================================================

Before arriving at the current pipeline, three fundamental difficulties had to be
addressed.

\subsection{State-Dependence of Reward Distributions}

The original zooming bandit algorithm (recovered as $H=1$ in \citet{cao2021}) assumes
a \emph{single} fixed reward distribution over actions. Zooming refines cells around
arms with high average reward, which is correct only when that distribution is
stationary across pulls. In an MDP, however, two different states $x \ne x'$ may have
completely different optimal actions. If states are conflated, the adaptive
discretization focuses on regions that are near-optimal on average over all states, but
potentially highly suboptimal for any particular one. Formally, the gap function
$\mathrm{gap}_h(x,a) = V_h^*(x) - Q_h^*(x,a)$ depends on $x$; there is no
state-independent notion of a ``good'' action. The near-optimal set
$P^{Q^*}_{h,r}$ in Definition~\ref{def:zooming} is defined over the \emph{joint}
space $\mathcal{S} \times \mathcal{A}$ precisely because of this state-dependence.
Neglecting it yields an ill-defined near-optimal set and invalidates the zooming
dimension analysis of Corollary~\ref{cor:zooming_regret}.

\subsection{Curse of Dimensionality in Joint State--Action Zooming}

The algorithm of \citet{cao2021} operates on the joint space $\mathcal{S} \times
\mathcal{A}$ with a single partition. At each split, a ball $B$ of radius $r$ is
replaced by children of radius $r/2$ forming an $r/2$-net of $B$. The number of
children is controlled by the doubling constant $\Lambda(\mathcal{S}\times\mathcal{A})$
(Definition~1 of \citet{cao2021}), which for Euclidean $\mathbb{R}^d$ is
$\Lambda = \mathcal{O}(3^d)$. With a 25-dimensional kinematic observation and a
1-dimensional steering action, $d = 26$, so each split produces up to
$\mathcal{O}(3^{26}) \approx 2.5\times10^{12}$ children---completely intractable. The
regret bound of Theorem~1 in \citet{cao2021} itself depends polynomially on $\Lambda$,
so large joint dimension destroys the theoretical guarantee as well.

\subsection{Absence of Optimism in Policy-Gradient Methods}

The zooming algorithm requires that the action-value estimate be
\emph{optimistic}---an upper confidence bound on $Q^*$---so that greedy action
selection drives exploration toward under-visited, potentially high-value regions.
The Q-update in Algorithm~1 of \citet{cao2021} adds an explicit exploration bonus
$b_t \propto 1/\sqrt{t}$, ensuring the estimate dominates $Q^*$ with high probability.
Early experiments used Proximal Policy Optimization (PPO) with Generalized Advantage
Estimation (GAE). PPO computes advantage estimates $A^\pi$ under the \emph{current}
policy, which track the signed deviation from a baseline rather than an optimistic
ceiling on $Q^*$. Consequently, a greedy policy over PPO advantages has no formal
guarantee of visiting under-explored regions, violating the fundamental assumption
of the zooming analysis.

% ============================================================
\section{Our Algorithm}
% ============================================================

We now describe our three-phase \emph{Bisimulation-Guided Clustered Zooming SAC}
(BisimZoom) pipeline.

\subsection{Phase 1: Learning a Bisimulation Representation}

We train an encoder $\varphi_\theta : \mathcal{S} \to \mathcal{Z}$ (with
$\mathcal{Z} = \mathbb{R}^d$, $d \ll \dim(\mathcal{S})$) jointly with a Soft
Actor--Critic (SAC) policy. Following Equation~(4) of \citet{zhang2021}, the encoder
loss on a batch of state pairs $(s_i, s_j)$ is
\[
  \mathcal{J}(\varphi) \;=\;
    \mathbb{E}_{s_i, s_j}\!\left[
      \Bigl(
        \|\varphi(s_i) - \varphi(s_j)\|_1
        - |r_i - r_j|
        - \gamma\,W_2\!\left(\hat{P}(\cdot|\bar{z}_i, a_i),\,
                            \hat{P}(\cdot|\bar{z}_j, a_j)\right)
      \Bigr)^{\!2}
    \right],
\]
where $\hat{P}$ is a learned Gaussian dynamics model and $\bar{z}$ denotes a
stop-gradient copy. The $W_2$ term admits a closed form for Gaussians:
$W_2(\mathcal{N}(\mu_i,\Sigma_i), \mathcal{N}(\mu_j,\Sigma_j))^2 =
\|\mu_i-\mu_j\|_2^2 + \|\Sigma_i^{1/2}-\Sigma_j^{1/2}\|_F^2$.
Theorem~1 of \citet{zhang2021} guarantees that $\varphi$ converges to a fixed-point
encoder satisfying $\|\varphi(s_i)-\varphi(s_j)\|_1 = \tilde{d}(s_i, s_j)$, provided
the policy is continuously improving.

\subsection{Phase 2: Clustering the Latent Space}

After Phase 1 we freeze $\varphi$ and embed the replay buffer:
$Z = \{\varphi(s) : s \in \mathcal{D}\}$.
We run $k$-means on $Z$ to obtain $k$ cluster centers $\{c_1,\ldots,c_k\}$ and
assign each state a cluster label. The theoretical justification is the value bound
of Theorem~2 of \citet{zhang2021}: aggregating states within an $\varepsilon$-ball
of the bisimulation metric incurs optimal-value error at most
\begin{equation}
  |V^*(s) - V^*(\varphi(s))| \;\le\; \frac{2\varepsilon}{(1-\gamma)(1-c)}.
  \label{eq:value_bound}
\end{equation}
Choosing $\varepsilon$ equal to the intra-cluster bisimulation diameter bounds the
accuracy of the resulting per-cluster policies. We select $k$ via the silhouette score
over the embedded buffer.

\subsection{Phase 3: Clustered Zooming SAC}

For each cluster $c \in [k]$ we maintain: (i) an \texttt{ActionZooming} instance with
active cubes $\{B_j^{(c)}\}$, visit counts $n_j^{(c)}$, and Q-estimates; (ii) twin
Q-networks and a categorical actor with rebuildable output heads; (iii) a per-cluster
replay buffer storing continuous actions, preserving validity after cube splits; and
(iv) an auto-tuned entropy temperature $\alpha^{(c)}$.

The critic is updated \emph{pessimistically} (minimum of twins) for stability. The
actor is updated \emph{optimistically}: a UCB bonus
\[
  u_j^{(c)} \;=\; \beta\sqrt{\frac{\log N^{(c)}}{n_j^{(c)}}},
  \qquad N^{(c)} = \textstyle\sum_j n_j^{(c)},
\]
is added to $Q$ before computing the actor loss. This directly encodes the optimism
principle of the zooming analysis and corresponds to the $b_t$ bonus in Algorithm~1
of \citet{cao2021}. The split threshold $N_{\mathrm{split}}(B) = (d_{\max}/r(B))^2$
mirrors the criterion of \citet{cao2021}, ensuring a cell is refined only after
enough samples have been collected to trust the local Q-estimate.

\begin{algorithm}[H]
\caption{BisimZoom: Bisimulation-Guided Clustered Zooming SAC}
\begin{algorithmic}[1]
\Require Environment, feature dim $d$, cluster count $k$, UCB coefficient $\beta$
\vspace{0.4em}
\Statex \textbf{Phase 1 --- Learn bisimulation encoder}
\State Train BisimSAC for $T_1$ steps; obtain frozen encoder $\varphi$
\vspace{0.3em}
\Statex \textbf{Phase 2 --- Cluster latent space}
\State Embed replay buffer $Z \gets \{\varphi(s) : s \in \mathcal{D}\}$
\State Run $k$-means on $Z$; store centers $\{c_1,\ldots,c_k\}$ and labels
\vspace{0.3em}
\Statex \textbf{Phase 3 --- Clustered Zooming SAC}
\For{$c = 1$ \textbf{to} $k$}
  \State Initialize \texttt{ActionZooming}$^{(c)}$, twin Q-nets, actor, replay buffer
\EndFor
\For{step $t = 1$ \textbf{to} $T_3$}
  \State Observe $s_t$; compute $z_t = \varphi(s_t)$;
         assign $c_t = \arg\min_c \|z_t - c_c\|$
  \If{$|\mathcal{D}^{(c_t)}| < N_{\min}$}
    \State $j \gets \mathrm{Uniform}\{1,\ldots,|\{B_j^{(c_t)}\}|\}$
      \Comment{random warm-up}
  \Else
    \State $j \gets \arg\max_j\; \pi^{(c_t)}(j \mid z_t)$
  \EndIf
  \State $a_t \gets \mathrm{center}(B_j^{(c_t)})$;\quad
         $n_j^{(c_t)} \mathrel{+}= 1$
  \State Execute $a_t$; observe $r_t, s_{t+1}$; compute $z_{t+1}, c_{t+1}$
  \State Store $(z_t, a_t, r_t, z_{t+1}, c_{t+1})$ in $\mathcal{D}^{(c_t)}$
  \For{\textbf{each} cluster $c$ with $|\mathcal{D}^{(c)}| \ge N_{\min}$}
    \State \textbf{Critic} (pessimistic):
      $\mathcal{L}_Q = \mathrm{MSE}\!\left(Q_i^{(c)}(z,j),\;
        r + \gamma V^{(c')}(z')\right)$
    \State \textbf{Actor} (optimistic):
      $\mathcal{L}_\pi =
        \mathbb{E}_j\!\left[\alpha^{(c)} \log \pi^{(c)}(j|z)
          - \left(Q_{\min}^{(c)}(z,j) + u_j^{(c)}\right)\right]$
    \State Update $\alpha^{(c)}$ toward target entropy;\quad
           soft-update target networks
  \EndFor
  \State \textbf{Split check}: for each $c$, if $n_j^{(c)} \ge N_{\mathrm{split}}$,
         split $B_j^{(c)}$; rebuild heads, inherit parent weights with small noise
\EndFor
\end{algorithmic}
\end{algorithm}


% ============================================================
\section{Current Problems and Potential Directions}
% ============================================================

\subsection{Rare-State Neglect Under Class Imbalance}

$k$-means minimizes intra-cluster $\ell_2$ variance, so clusters grow to enclose
densely populated regions of $\mathcal{Z}$. Rare but safety-critical states---a
vehicle suddenly veering into the ego lane---contribute negligibly to the objective
and are either absorbed into a large common cluster or assigned an undersized cluster
with insufficient replay data. The value bound \eqref{eq:value_bound} still holds
theoretically, but the per-cluster policy for the rare state will be poorly trained
due to sample starvation: the bound is vacuous in practice when the cluster receives
almost no gradient updates.

A related observation from \citet{zhang2021} is that the bisimulation encoder is
trained to be invariant to task-irrelevant variation. This is desirable in general,
but may cause rare states that differ from common ones primarily through
task-\emph{relevant} features (e.g.\ a nearby vehicle's trajectory) to be
inappropriately compressed together with safer states if those features produce only
modest reward differences at the current policy.

\paragraph{Potential directions.}
\begin{itemize}[noitemsep]
  \item \textbf{Risk-sensitive clustering.} Replace $k$-means with a
    density-weighted objective, such as $k$-medoids with an exponential cost
    on bisimulation-distance to states with high reward variance, explicitly
    preventing high-stakes rare states from being subsumed by large common clusters.
  \item \textbf{Prioritized replay with importance weighting.} Adopt prioritized
    experience replay within each cluster buffer, weighting samples by TD-error or
    by the inverse empirical visit frequency of the corresponding cluster. This
    ensures that transitions from rare clusters receive proportionally more gradient
    signal.
  \item \textbf{Forced splitting of high-variance clusters.} After standard
    clustering, run a secondary pass identifying clusters with high reward
    variance and force-split them, analogous to the buffering phase of
    \citet{cao2021} where children are initialized with large bonuses before being
    promoted to active use.
\end{itemize}

\subsection{Sensitivity to Hyperparameters $d$ and $k$}

The current pipeline requires manual selection of the latent dimension $d$ and the
number of clusters $k$. An observed failure mode is that over 95\% of the variance in
$\varphi(s)$ is explained by 1--2 principal components, suggesting the encoder
collapses to a near-one-dimensional representation.

This observation has an important positive implication. If the effective rank of
$\varphi(\mathcal{D})$ is $d^* \ll d$, then the true intrinsic dimension of the
latent space is small, and with a 1-dimensional action space the effective joint
dimension for zooming is $d^* + 1$. By Corollary~\ref{cor:zooming_regret}, the
regret depends on the zooming dimension $z_c \le d^* + 1$. For $d^* = 2$ this gives
a joint dimension of 3, with doubling constant $\Lambda = \mathcal{O}(3^3) = 27$
rather than $\mathcal{O}(3^{26})$. In this low-dimensional regime, the direct
joint-space Adaptive Q-learning of \citet{cao2021} becomes fully tractable, and
there would be no need for the clustering approximation at all.

\begin{remark}
The practical viability of direct joint zooming hinges on whether $d^*$ is genuinely
small or whether the 95\% threshold obscures important rare-state variation that lives
in the remaining 5\%. This is exactly the tension between Problems~1 and~2: the
dimensions discarded as low-variance may be precisely those encoding the rare but
critical states discussed above.
\end{remark}

\paragraph{Potential directions.}
\begin{itemize}[noitemsep]
  \item \textbf{Automatic dimensionality selection.} Compute the eigenspectrum of the
    covariance of $\varphi(\mathcal{D})$ and set
    $d^* = \min\{d' : \sum_{i\le d'}\lambda_i / \sum_i \lambda_i \ge 1-\delta\}$
    for a threshold $\delta$ chosen to preserve rare-state variation (e.g.\
    $\delta = 0.01$), then compress $\mathcal{Z}$ to $\mathbb{R}^{d^*}$ before
    clustering.
  \item \textbf{Direct joint zooming when $d^*$ is small.} If $d^* + \dim(\mathcal{A})
    \le 4$ or $5$, replace clustered independent zooming with Adaptive Q-learning
    \citep{cao2021} on the compressed joint space, recovering the full
    zooming-dimension guarantee without the clustering approximation.
  \item \textbf{Information-theoretic cluster selection.} Minimize a penalized
    within-cluster bisimulation-diameter criterion instead of Euclidean silhouette
    score, grounding the choice of $k$ in the value bound \eqref{eq:value_bound}
    of \citet{zhang2021}.
\end{itemize}

\subsection{Representation--Policy Gradient Conflict in BisimSAC}

In Phase 1, plain SAC achieves higher reward than BisimSAC, despite the only
difference being the addition of the bisimulation encoder loss. We identify three
plausible causes.

First, \textbf{gradient conflict}: the encoder $\varphi$ receives simultaneous
gradients from the bisimulation loss $\mathcal{J}(\varphi)$ and the SAC critic
(Algorithm~2 of \citet{zhang2021} allows critic gradients to flow back through
$\varphi$). The bisimulation loss pushes $\varphi$ to compress task-irrelevant
variation, while the critic may rely on fine-grained state distinctions for precise
value estimation. These objectives can directly oppose each other.

Second, \textbf{noisy $W_2$ targets}: the dynamics model $\hat{P}$ must be accurate
for the bisimulation loss to provide a useful signal. Early in training, $\hat{P}$ is
unreliable, so the $W_2$ term introduces large variance into $\mathcal{J}(\varphi)$
and corrupts the encoder before the policy has had time to stabilize.

Third, \textbf{premature compression}: on a 25-dimensional kinematic observation
where all features are task-relevant, the bisimulation loss may compress states that
differ only in a distant vehicle's current lane position into the same latent point,
even though that vehicle may become relevant within a few timesteps. This degrades
the critic's ability to discriminate near-future outcomes.

\paragraph{Potential directions.}
\begin{itemize}[noitemsep]
  \item \textbf{Stop-gradient encoder.} Train the critic on $\bar{\varphi}(s)$
    (stop-gradient copy) and update $\varphi$ only through $\mathcal{J}(\varphi)$,
    eliminating the gradient conflict. This is explicitly noted as optional in
    \citet{zhang2021}: backpropagating through $\varphi$ ``can improve performance
    further'' but is not always beneficial.
  \item \textbf{Delayed and annealed encoder training.} Begin encoder training only
    after the policy and dynamics model have stabilized (e.g.\ after $10^4$ steps),
    so that $\hat{P}$ provides a meaningful $W_2$ target. Linearly anneal the
    coefficient on $\mathcal{J}(\varphi)$ from $0$ to $1$ over training to prevent
    early corruption of the encoder.
  \item \textbf{Separate encoder pre-training.} Pre-train $\varphi$ offline on a
    diverse dataset before beginning policy optimization, so that the bisimulation
    representation is already stable when SAC training starts, removing the
    co-adaptation problem entirely.
\end{itemize}



\newpage

\bibliography{references}
\bibliographystyle{abbrvnat}

\end{document}
hi taro!
:)
